\documentclass[11pt,a4paper]{jsarticle}

\usepackage[top=28mm,bottom=28mm,left=25mm,right=25mm]{geometry}

% 数式・定理系を先に読む
\usepackage{amsmath}
\usepackage{amssymb}
\usepackage{amsthm}

% hyperref 系はその後
\usepackage[dvipdfmx]{hyperref}
\usepackage{pxjahyper}

% cleveref は hyperref と amsmath の後
\usepackage{cleveref}

\hypersetup{
  colorlinks=true,
  linkcolor=blue,
  citecolor=blue,
  urlcolor=blue,
  pdftitle={Follmer過程によるDDPMの内在次元適応的Wasserstein収束解析},
  pdfauthor={Anonymous}
}
\newtheorem{definition}{定義}[section]
\newtheorem{theorem}{定理}[section]
\newtheorem{proposition}[theorem]{命題}
\newtheorem{lemma}[theorem]{補題}
\newtheorem{corollary}[theorem]{系}
\newtheorem{assumption}[theorem]{仮定}
\theoremstyle{remark}
\newtheorem{remark}[theorem]{注意}
\newtheorem{example}[theorem]{例}
\renewcommand{\proofname}{証明}
\renewcommand{\abstractname}{概要}
\crefname{theorem}{定理}{定理}
\Crefname{theorem}{定理}{定理}
\crefname{lemma}{補題}{補題}
\Crefname{lemma}{補題}{補題}
\crefname{corollary}{系}{系}
\Crefname{corollary}{系}{系}
\crefname{assumption}{仮定}{仮定}
\Crefname{assumption}{仮定}{仮定}
\crefname{remark}{注意}{注意}
\Crefname{remark}{注意}{注意}
\crefname{equation}{式}{式}
\Crefname{equation}{式}{式}

\DeclareMathOperator{\Tr}{Tr}
\DeclareMathOperator{\Cov}{Cov}
\DeclareMathOperator{\Law}{Law}
\DeclareMathOperator{\KL}{KL}
\DeclareMathOperator{\supp}{supp}
\DeclareMathOperator{\diam}{diam}
\newcommand{\R}{\mathbb{R}}
\newcommand{\E}{\mathbb{E}}
\newcommand{\Pbb}{\mathbb{P}}
\newcommand{\cN}{\mathcal{N}}
\newcommand{\cP}{\mathcal{P}}
\newcommand{\cX}{\mathcal{X}}
\newcommand{\cF}{\mathcal{F}}
\newcommand{\Wtwo}{W_2}
\newcommand{\dd}{\mathrm{d}}
\newcommand{\eps}{\varepsilon}
\newcommand{\wh}{\widehat}

\title{内在次元を用いた拡散モデルのWasserstein収束解析}
\author{枇榔一真}
\date{2026年6月20日}

\begin{document}
\maketitle
\section{設定}
\subsection{SDEを通じたDDPMの解釈}
まず, 通常のDDPMの更新則を設定する. 
\begin{equation}
    X_t\stackrel{\mathrm{d}}{=}\sqrt{1-\sigma^2_t}X_0+\sigma_tW \quad (\sigma_t=\sqrt{1-e^{-2t}})
\end{equation}
ここで, $X_0 \sim p_{data}, W\sim \mathcal{N}(0,I_d)$ とする.
これは, データをノイズ化していく過程とみなせ, $X_t$ の分布は $t$ が大きくなるにつれ $\mathcal{N}(0,I_d)$ に近づく.
次に, $0=t_0<t_1<\dots < t_N<t_{N+1}=T$ とする.
\begin{align}
    &Y_{t_0}\sim \mathcal{N}(0,I_d), \\
    &Y_{t_n+1}=\frac{1}{\sqrt{\alpha_n}}(Y_{t_n}+(1-\alpha_n)\hat{s}_{T-t_n}(Y_{t_n}))+\sqrt{\frac{(1-\alpha_n)(1-\bar{\alpha}_{n+1})}{1-\bar{\alpha_n}}}Z_n\quad(n=1,\dots, N)
\end{align}
ここで, $\hat{s}_t$ は スコア関数 $s_t=\nabla \log{p_{t}}$ とする. また, $p_t$ は $X_t$ の密度とする.
このとき, $Y$ はデータを生成する過程と考えられ, 
以降, 
\begin{equation}
    \alpha_n=e^{-2(t_{n+1}-t_n)}, \bar{\alpha}_n=\Pi_{i=n}^N \alpha_i
\end{equation}
とする.


次に, このアルゴリズムのSDE的解釈を与える.
\begin{equation}
    dX_t=-X_tdt+\sqrt{2}dW_t, \quad X_0\sim p_{data}, \quad t\in[0,T]
\end{equation}
ここで, このSDEの弱解は(1)式で与えられる確率過程と, 各時刻で分布の意味で等しくなる.
次に, データを生成するプロセスに対応するSDEを次の式で与える.
\begin{equation}
    dY_t=(Y_t+2s_{T-t}(Y_t))dt+\sqrt{2}dB_t \quad t\in[0,T]
\end{equation}
この時, $Y_0\stackrel{\mathrm{d}}{=} X_T$ とすると, 任意の時刻 $t\in [0,T]$ において
$Y_{T-t}\stackrel{\mathrm{d}}{=}X_t$ が成立する(\cite{anderson1982reverse}, \cite{haussmann1986time}).
実際に (6) 式に基づいてデータを生成する際には, スコア関数を推定量 $\hat{s}_{T-t}$ に置き換え, 時間に関して離散化をする必要がある.

\subsection{スコア関数とTweedieの公式}
\begin{equation}
    \mu_t(x)=\mathbb{E}[X_0|X_t=x], \quad \Sigma_t(x)=Cov[X_0|X_t=x]
\end{equation}
とする. このとき, Tweedieの公式(\cite{efron2011})により, 
\begin{equation}
    \mu_t(X_t)=\frac{1}{\sqrt{1-\sigma^2_t}}(X_t+\sigma^2_t\nabla \log{q_t(X_t)})=\frac{1}{\sqrt{1-\sigma^2_t}}(X_t+\sigma^2_ts_t(X_t))
\end{equation}
が得られる.
\begin{equation}
    \hat{\mu}_t(X_t)=\frac{1}{\sqrt{1-\sigma^2_t}}(X_t+\sigma^2_t\hat{s}_t(X_t))
\end{equation}
これより, スコア関数は次のように書き換えられる.
\begin{equation}
    s_t(X_t)=\frac{\sqrt{1-\sigma^2_t}}{\sigma^2_t}\mu_t(X_t)-\frac{1}{\sigma^2_t}X_t
\end{equation}
よって, (6) 式は次のように書き換えられる.
\begin{align}
    dY_t&=(Y_t+\frac{2\sqrt{1-\sigma^2_{T-t}}}{\sigma^2_{T-t}}\mu_{T-t}(Y_t)-\frac{2}{\sigma^2_{T-t}}Y_t)dt+\sqrt{2}dB_t\\
    &=\{(1-\frac{2}{\sigma^2_{T-t}})Y_t+\frac{2\sqrt{1-\sigma^2_{T-t}}}{\sigma^2_{T-t}}\mu_{T-t}(Y_t)\}dt+\sqrt{2}dB_t\\
\end{align}

\subsection{集合の内在次元について}
データ分布の内在次元について議論するために, $p_{data}$ のサポート集合の複雑さを定義する.
\begin{definition}
    $A\subset \mathbb{R}^d$ に対して, スケール $\epsilon_0>0$ における convering number を$N^{cover}(A,\|\|_2,\epsilon_0)$ とかく.
    これを, 
    \begin{equation}
        A\subset \cup_{i=1}^sB(x_i,\epsilon_0)
    \end{equation}
    を満たすような $x_1,\cdots, x_s$ が存在するような最小の $s$ として定義する.
\end{definition}
上記のcovering numberに関して, 次の仮定をおく.
\begin{assumption}
    
\end{assumption}

\begin{assumption}

\end{assumption}


\section{収束解析}
推定スコア $\hat{s}_t$ に対して次のリプシッツ連続性の仮定をおく.
\begin{assumption}
    ある二乗可積分な関数 $L:(0,1)\rightarrow [0,\infty)$ が存在して, 次の式が成立する.
    \begin{equation}
        \|\hat{s}_t(x)-\hat{s}_t(y)\|\leqq L(t)\|x-y\| \quad  ^{\forall}(t,x,y)\in (0,1)\times \mathbb{R}^d\times \mathbb{R}^d
    \end{equation}
\end{assumption}
また, スコア推定に関する次の仮定を置く.
\begin{assumption}
    \begin{equation}
        \sum_{n=0}^{N-1}(t_{n+1}-t_n)\mathbb{E}[\|s_{T-t_n}(X_{T-t_n})-\hat{s}_{T-t_n}(X_{T-t_n})\|^2]\leqq \epsilon^2<\infty
    \end{equation}
\end{assumption}
これらの仮定のもとで, 次の定理を証明する.
\begin{theorem}
    \begin{equation}
        W_2(p_{data}, Y_{t_N})\leqq 
    \end{equation}
\end{theorem}
\begin{proof}
    
\end{proof}
\end{document}







\begin{thebibliography}{99}

\bibitem{efron2011}
Efron, B.
\newblock Tweedie’s formula and selection bias.
\newblock Journal of the American Statistical Association,
106(496):1602–1614.

\bibitem{huang2026ddpm}
Zhihan Huang, Yuting Wei, and Yuxin Chen.
\newblock Denoising diffusion probabilistic models are optimally adaptive to unknown low dimensionality.
\newblock arXiv preprint arXiv:2410.18784, 2026.

\bibitem{ho2020ddpm}
Jonathan Ho, Ajay Jain, and Pieter Abbeel.
\newblock Denoising diffusion probabilistic models.
\newblock Advances in Neural Information Processing Systems, 33:6840--6851, 2020.

\bibitem{song2021score}
Yang Song, Jascha Sohl-Dickstein, Diederik P. Kingma, Abhishek Kumar,
Stefano Ermon, and Ben Poole.
\newblock Score-based generative modeling through stochastic differential equations.
\newblock International Conference on Learning Representations, 2021.

\bibitem{anderson1982reverse}
Brian D. O. Anderson.
\newblock Reverse-time diffusion equation models.
\newblock Stochastic Processes and their Applications, 12(3):313--326, 1982.

\bibitem{haussmann1986time}
Ulrich G. Haussmann and Etienne Pardoux.
\newblock Time reversal of diffusions.
\newblock The Annals of Probability, 14(4):1188--1205, 1986.

\bibitem{oksendal2003sde}
Bernt {\O}ksendal.
\newblock Stochastic Differential Equations: An Introduction with Applications.
\newblock Springer, 2003.

\end{thebibliography}
